% Part: first-order-logic 
% Chapter: axiomatic-deduction 
% Section: proving-things-quant

\documentclass[../../../include/open-logic-section]{subfiles}

\begin{document}

\olfileid{fol}{axd}{prq}

\olsection{\usetoken{P}{derivation} mawa Kuantor}

\begin{ex}
Ayo kita menehi !!a{derivation} saka $(\lforall[x][!A(x)] \land
\lforall[y][!B(y)]) \lif \lforall[x][(!A(x) \land !B(x))]$.

Kaping pisan, gatekna yen
\begin{align*}
  (\lforall[x][!A(x)] \land \lforall[y][!B(y)]) & \lif \lforall[x][!A(x)]\\
  \intertext{iku conto saka \olref[prp]{ax:land1}, lan}
  \lforall[x][!A(x)] & \lif !A(a) \\
  \intertext{conto saka \olref[qua]{ax:q1}. Mula, miturut \olref[pro]{prop:chain}, kita ngerti yen}
  (\lforall[x][!A(x)] \land \lforall[y][!B(y)]) & \lif !A(a) \\
  \intertext{bisa !!{derivable}. Semono uga, amarga}
  (\lforall[x][!A(x)] \land \lforall[y][!B(y)]) & \lif \lforall[y][!B(y)] \qquad\text{lan}\\
    \lforall[y][!B(y)] & \lif !B(a)\\
    \intertext{iku conto saka \olref[prp]{ax:land2} lan \olref[qua]{ax:q1}, miturut urutane,}
    (\lforall[x][!A(x)] \land \lforall[y][!B(y)]) & \lif !B(a)\\
    \intertext{bisa diderivasi miturut \olref[pro]{prop:chain}. Kanthi nggunakake conto kang trep saka \olref[prp]{ax:land3} lan rong aplikasi~\MP, kita weruh yen}
    (\lforall[x][!A(x)] \land \lforall[y][!B(y)]) & \lif (!A(a) \land !B(a))\\
    \intertext{bisa diderivasi. Saiki kita bisa ngetrapake \QR{} kanggo entuk}
(\lforall[x][!A(x)] \land \lforall[y][!B(y)]) & \lif \lforall[x][(!A(x) \land !B(x))].
\end{align*}
\end{ex}

\end{document}
